\section{The penalty threshold scales with \texorpdfstring{$\rho$}{rho}}\label{sec:penalty}

\begin{proposition}\label{prop:penalty}\leavevmode
\begin{enumerate}[label=(\alph*)]
\item \emph{Sufficiency.} $\mu\ge\mu_0=4\kappa C_I^2\rho$ implies coercivity (Lemma~\ref{lem:coercive}).
\item \emph{Necessity, with an exact certificate.} Let $z$ be a vertex of a straight part of the Nitsche boundary, and let $\ell_z$ be the length of its boundary edges. Suppose its star, dilated by the factor $1/\ell_z$, is congruent to one of the reference stars $\cS_3$ or $\cS_4'$ below. If $\mu/(h/\ell_z)<9$, then $N_h$ is indefinite on $\Zh$.

The same holds, by continuity, at a vertex of $\Gh$ with interior angle $\pi+O(\hG)$ whose star is a small perturbation of $\cS_3$ or $\cS_4'$. This uses that both reference stars are non-singular: every vertex, including the rim vertices, has non-collinear incident edges. The dimension of the divergence-free subspace therefore equals the generic count and is locally constant. This fails at the singular $m=2$ configuration, which (M1$'$) excludes.
\end{enumerate}
\end{proposition}

So, for meshes that contain such a boundary star at a vertex with $\rho\asymp h/\ell_z$, the condition $\mu_0\asymp\rho$ is both necessary and sufficient. Theorem~\ref{pb:thm} and Corollary~\ref{pb:rho} below replace the star hypothesis by a condition on a single boundary triangle. The threshold is local: what is required is $\mu/h\gtrsim c/\ell_z$. The $\rho$-scaling is the price of Scott's normalisation of the penalty by the global $h$.

\begin{proof}[Proof of (b)]
\begin{enumerate}
\item In two dimensions, $a(v,v)$ and $\ip{\dn v}{v}$ are invariant under dilation, and $\norm{v}^2_{L^2(\Gh)}$ scales linearly.
\item A witness $v$ on the reference star, extended by zero, lies in $\Zh$: it is continuous, vanishes on the star's outer chords and on $\partial R$, and is divergence-free.
\item It satisfies $N_h(v,v)=a(v,v)-2\ip{\dn v}{v}+(\mu\ell_z/h)\norm{v}^2_{\mathrm{ref}}$.
\item By the certificate below, $2B-A-9C>0$ on each reference star, so $N_h(v,v)<0$ when $\mu\ell_z/h<9$.\qedhere
\end{enumerate}
\end{proof}

\subsection*{The exact certificate}
The reference stars are fans at $O=(0,0)$ with rational rim points:
\begin{itemize}
\item $\cS_3$: rim points $(1,0),(1,1),(-1,1),(-1,0)$;
\item $\cS_4'$: rim points $(1,0),(1,1),(0,\tfrac65),(-1,1),(-1,0)$.
\end{itemize}
The apex of $\cS_4'$ is moved off the line $y=1$. With apex $(0,1)$ the rim vertex is singular, the kernel dimension jumps from $16$ to $17$, and the perturbation argument does not apply.

The Nitsche boundary is $\{y=0\}$, and the rim chords carry zero Dirichlet data. The velocity is represented on each triangle by $P_4$ polynomials in global coordinates with rational coefficients. The computation is done in exact rational arithmetic (SymPy):
\begin{itemize}
\item Continuity across the spokes, vanishing on the rim, and pointwise $\div v=0$ are imposed as polynomial identities. This gives $80$ and $110$ linear constraints and kernel dimensions $14$ and $16$, the generic counts.
\item The quadratic forms $A=\sum_T\int\frac12Dv:Dv$, $B=\int(\dn v)\cdot v$ and $C=\int|v|^2$ are integrated exactly.
\item A rational witness in the kernel satisfies $2B-A-9C>0$ \emph{exactly} on each star. Its exact Rayleigh quotients $(2B-A)/C$ are $9.97854\ldots$ ($\cS_3$) and $9.26747\ldots$ ($\cS_4'$); these are lower bounds for the largest Rayleigh quotient $\lambda^\ast$ of each star.
\end{itemize}
\cert{code/lamstar\_exact.py, code/lamstar\_exact\_v2.py}{logs/lamstar\_exact\_*.log}
The witnesses are continuous piecewise-$P_4$ and divergence-free, so they lie in $\Zh$ for every $k\ge4$. The necessity statement therefore holds for all $k\ge4$; for $k>4$ the true local threshold is at least as large.

In floating point, the largest Rayleigh quotient $\lambda^\ast$ is about $15$ on uniform fans with $m=3,4$ and ranges over $14.8$--$38.5$ on random $4$-fans (\texttt{code/lamstar.py}). As a global check, the measured threshold $\gamma^\ast=\mu^\ast\hG/h$ stays in $18.1$--$20.1$ as $h/\hG$ varies over $3.3$--$44$ on two mesh levels. Under refinement at fixed $h/\hG\approx3.3$ it lies in $14.7$--$20.8$, with increments that roughly halve per doubling (Table~\ref{tab:thresh}). Section~\ref{pb:sec} removes the reference-star hypothesis and explains the gap between the certified $9$ and the measured $18$--$21$.

\subsection{Penalty necessity beyond two reference stars}\label{pb:sec}

Proposition~\ref{prop:penalty}(b) certifies indefiniteness of $N_h$ on $\Zh$ for $\mu\ell_z/h<9$ only at vertices whose star is
congruent to (or a perturbation of) $\cS_3$ or $\cS_4'$. We replace it by a statement about \emph{every} boundary triangle, with
an explicit closed-form constant and an exact certificate over a two-parameter family of shapes. We then explain the gap
between the local constant and the global $\gamma^\ast\approx18$--$21$ of Table~\ref{tab:thresh}.

\subsubsection{A witness supported in one boundary triangle}

Let $T\in\Th$ have an edge $e\subset\Gh$, and let $T$ not meet $\partial R$. Rotate, translate and dilate by $1/|e|$ so that
$e=[(0,0),(1,0)]$, the fluid lies above and the apex is $(a,b)$, $b>0$. Let $\lambda_0,\lambda_1,\lambda_2$ be the barycentric
coordinates of the apex and of the base vertices $(0,0)$, $(1,0)$. For $c\in\R^3$ put
\[
\psi_c:=\lambda_1^2\lambda_2^2\,(c_0\lambda_0+c_1\lambda_1+c_2\lambda_2)\in P_5,\qquad v_c:=\curl\psi_c\in P_4^2 .
\]
Since $\psi_c$ vanishes to second order on the two interior edges $\{\lambda_1=0\}$ and $\{\lambda_2=0\}$, $v_c$ vanishes there.
Extended by zero, $v_c$ is continuous, zero on $\partial R$ and pointwise divergence-free. So $v_c\in\Zh$ for every $k\ge4$ and
every mesh containing $T$. No condition on the vertices (singular or not), on the other triangles, or on the straightness of
$\Gh$ beyond the edge $e$ itself is needed.

Let $A_T=\int_T\frac12Dv:Dv$, $B_T=\int_e(\dn v)\cdot v$ and $C_T=\int_e|v|^2$, with $n=(0,-1)$, as quadratic forms in $c$. They are
computed exactly (SymPy) and, writing $\varpi:=4a^2+4b^2$, are
\begin{align*}
C_T&=\frac1{1260\,b^2}\begin{pmatrix}2&a-6&-(a+5)\\ a-6&2(\varpi-3a+9)&\varpi-4a+15\\ -(a+5)&\varpi-4a+15&2(\varpi-5a+10)\end{pmatrix},\\
B_T&=\frac1{840\,b^3}\begin{pmatrix}12&-(\varpi-6a+27)&-(\varpi-2a+25)\\ -(\varpi-6a+27)&6(\varpi-3a+9)&3(\varpi-4a+15)\\ -(\varpi-2a+25)&3(\varpi-4a+15)&6(\varpi-5a+10)\end{pmatrix},
\end{align*}
and $A_T=(\alpha_{ij})/b^3$ with
\begin{align*}
630\,\alpha_{00}&=3a^4-6a^3+6a^2b^2+12a^2-6ab^2-9a+3b^4+4b^2+9,\\
420\,\alpha_{01}&=3a^4+6a^2b^2-5a^2+6a+3b^4-7b^2-9,\\
420\,\alpha_{02}&=3a^4-12a^3+6a^2b^2+13a^2-12ab^2-8a+3b^4-b^2-5,\\
70\,\alpha_{11}&=3a^4-3a^3+6a^2b^2+5a^2-3ab^2-3a+3b^4+7b^2+3,\\
140\,\alpha_{12}&=a^4-2a^3+2a^2b^2+5a^2-2ab^2-4a+b^4+7b^2+5,\\
70\,\alpha_{22}&=3a^4-9a^3+6a^2b^2+14a^2-9ab^2-10a+3b^4+10b^2+5 .
\end{align*}
$\det C_T=(3a^4-6a^3+7a^2b^2+3a^2-7ab^2+4b^4+5b^2)/(83\,349\,000\,b^6)>0$. (The forms were checked against an independent
quadrature/finite-difference evaluation, relative difference $4.5\cdot10^{-9}$, and $\lambda_T$ below against our
penalty-threshold code on single-triangle patches of the Section~\ref{sec:numerics} meshes: $10.48912$ vs.\ $10.48912$.)

\begin{theorem}[Single-triangle necessity]\label{pb:thm}
Let $\lambda_T$ be the largest eigenvalue of the pencil $(2B_T-A_T,\,C_T)$; it depends only on the shape of $T$. If
$\mu|e|/h<\lambda_T$, then $N_h$ is indefinite on $\Zh$ for every $k\ge4$. Consequently every mesh satisfies the necessary
condition
\[
\mu\ \ge\ h\,\max_{e\in\EG}\frac{\lambda_{T_e}}{|e|} .
\]
\end{theorem}

\begin{proof}
In two dimensions $A$ and $B$ are dilation invariant and $C$ scales like $|e|$. So for the eigenvector $c$,
$N_h(v_c,v_c)=A_T-2B_T+(\mu|e|/h)C_T=(\mu|e|/h-\lambda_T)\,C_T(c)<0$.
\end{proof}

\begin{proposition}[Exact certificates over shape families]\label{pb:cert}
For $\varphi\in(0,\pi)$ let $Q_\varphi$ be the set of apex positions $(a,b)$ for which the apex angle is at least $\varphi$ and $b\ge
s\max(a,1-a)$, with $s=0.087<\tan5^\circ$. This contains every triangle with apex angle $\ge\varphi$ and base angles $\ge5^\circ$.
Then
\[
\lambda_T>9.3\ \text{on }Q_{60^\circ},\qquad \lambda_T>4\ \text{on }Q_{45^\circ},\qquad \lambda_T>1\ \text{on }Q_{35^\circ}.
\]
\end{proposition}

The certificate proves only the three strict lower bounds. As numerical observations, not certified (\texttt{logs/lamstar\_family\_levels.log}): the first bound is nearly sharp, since a scan gives the minimum $9.31381$ over $Q_{60^\circ}$, at the equilateral triangle; along the level sets of the apex angle the scanned minima are $6.07$ ($50^\circ$), $4.47$ ($45^\circ$), $2.83$ ($40^\circ$), $1.12$ ($35^\circ$) and $-0.73$ ($30^\circ$), each at the isosceles shape.

\begin{proof}[Proof (computer-assisted, exact)]
$Q_\varphi$ is the part of the disk through $(0,0)$, $(1,0)$ with inscribed angle $\varphi$ that lies above the two lines. Cover a
rational box containing $Q_\varphi$ by dyadic sub-boxes. A box is discarded only if it provably misses $Q_\varphi$, using
outward-rounded rational enclosures of the disk. On every other box, take a rational vector $q$ (the eigenvector at the box
centre, rounded to denominators $\le1000$) and verify
\[
\Phi_q(a,b):=2520\,b^3\,q^{\mathsf T}\bigl(2B_T-A_T-c\,C_T\bigr)q>0
\]
on the whole box. $\Phi_q$ is a polynomial with rational coefficients (the denominators clear exactly), and the lower bound is
evaluated monomial-wise in exact rational interval arithmetic. The boxes are split until this succeeds.
\begin{itemize}
\item $c=9$ and $c=9.3$ on $Q_{60^\circ}$: $323$ and $716$ boxes, depth $\le11$.
\item $c=4$ on $Q_{45^\circ}$: $259$ boxes.
\item $c=1$ on $Q_{35^\circ}$: $297$ boxes.
\end{itemize}
As negative controls, the same procedure fails for $c=9.4$ on $Q_{60^\circ}$ and for $c=4.5$ on $Q_{45^\circ}$, as it must.
\end{proof}
\cert{code/lamstar\_family.py bubble; code/lamstar\_family.py certify}{logs/lamstar\_family\_bubble.log, logs/lamstar\_family\_certify.log}

\begin{corollary}[$\rho$-scaling is necessary on a broad class]\label{pb:rho}
Let $e_\ast$ be a shortest edge of $\Gh$ and suppose its triangle $T_{e_\ast}$ has apex angle $\ge35^\circ$ and base angles $\ge5^\circ$, and does not meet $\partial R$ (automatic for small $h$, since $\diam T_{e_\ast}\le C\hG<L-1$).
Then coercivity of $N_h$ on $\Zh$ requires $\mu>\rho$. If the apex angle is $\ge45^\circ$ (resp.\ $\ge60^\circ$) it requires
$\mu>4\rho$ (resp.\ $\mu>9.3\rho$). Together with Lemma~\ref{lem:coercive}, $\mu_0\asymp\rho$ is necessary and sufficient on
every such mesh, with no star-congruence or perturbation hypothesis.
\end{corollary}

\begin{remark}[What the single triangle cannot do]\label{pb:limits}
$\lambda_T<0$ for tall triangles, for example every right triangle with the right angle at a base vertex and apex angle
$\le30^\circ$ (numerically $\lambda_T\approx-0.2500$ at apex $30^\circ$, decreasing to $-15.7$ at $10^\circ$). A uniform statement under (M0)--(M2) alone would need witnesses on
several triangles. In the Section~\ref{sec:numerics} meshes such tall boundary triangles occur near the corner directions;
at $N=64$ the smallest value is $\lambda_T=-3.8$, with scaled apex about $(-0.08,2.38)$ (\texttt{logs/lamstar\_family\_levels.log}). There the vertex stars and first-layer windows still
have positive $\lambda^\ast$. Earlier versions conjectured that $\lambda^\ast_{\text{star}}\ge c(\Theta_0)>0$ for every nonsingular boundary star. That conjecture is false (Proposition~\ref{pb:counter}).
\end{remark}

\subsubsection{Stars need not carry a witness}

\begin{proposition}[Exact counterexamples]\label{pb:counter}
Let $\cS^{(H)}$ be the three-triangle column star at $O=(0,0)$ with rim points $(1,0),(1,H),(0,H),(-1,0)$ and Nitsche boundary $\{y=0\}$; it is the star of a boundary vertex of a mesh of columns of width $1$ and height $H$ split by the forward diagonal. It satisfies (M1), (M1$'$) and (M2) with $\Theta=1$ at every vertex; its minimum angle is $\arctan(1/H)$. In exact rational arithmetic ($LDL^{\mathsf T}$ without pivoting, all pivots positive), on the divergence-free space of continuous piecewise-$P_k$ fields on the patch vanishing on the rim:
\begin{enumerate}[label=(\alph*)]
\item $k=4$: $A-2B-C\succ0$ on $\cS^{(5)}$ (minimum angle $11.3^\circ$) and $A-2B-\frac12C\succ0$ on $\cS^{(9/2)}$;
\item $k=4$, the union of the two stars of one boundary edge: $A-2B-\frac1{10}C\succ0$ at $H=6$ and $A-2B-C\succ0$ at $H=8$;
\item the star fails for $k=5$ at $H=8$ and $k=6$ at $H=12$; on Clough--Tocher split stars it fails for $k=2$ at $H=2$ (macro minimum angle $26.6^\circ$), $k=3$ at $H=4$ and $k=4$ at $H=8$.
\end{enumerate}
In each case no field supported in the patch is a witness for any $\mu\ge0$. Wider windows still work on the same meshes: three consecutive stars at $H=5$ give $2B-A-C>0$ and four stars at $H=8$ give $2B-A-\frac14C>0$, so $\mu\gtrsim\rho$ still holds there.
\end{proposition}

\begin{proof}
The forms and spaces are those of Section~\ref{sec:penalty} ($A=\sum_T\int\frac12Dv:Dv$, $B=\int(\dn v)\cdot v$ with $n=(0,-1)$, $C=\int|v|^2$), with continuity, vanishing on the rim and $\div v=0$ imposed as polynomial identities over $\mathbb Q$ and the null space computed exactly. The code reproduces the witness quotient $9.978$ of $\cS_3$, and negative controls ($-\frac65$ at $H=5$, $0$ on $\cS_3$ and on $\cS^{(4)}$, $-\frac3{20}$ for the $H=6$ edge patch) fail as they must.
\end{proof}
\cert{notes/v6/penalty/certs.py; notes/v6/penalty/exact.py}{notes/v6/penalty/cert\_all\_k.log, cert\_edge6.log, cert\_high.log}

The conclusions persist under small perturbations, including bending the boundary: for fixed combinatorics, all vertices of these patches are non-singular, so the dimension of the divergence-free space equals the generic count and the space depends continuously on the vertex positions (as in step~4 of Proposition~\ref{ct:prop:penalty}); the strict positivity $A-2B-cC\succ0$ is then an open condition. A global mesh of $\Oh$ containing such a star is obtained from a polar column layer of height ratio $H$ around $\Gh$; we have not written it out. The counterexamples refute $\lambda^\ast_{\text{star}}\ge c(\Theta_0)$ as stated, because the minimum angle there is small. They do \emph{not} refute a bound $c(\Theta_0,\theta_{\min})$ for $\theta_{\min}\gtrsim20^\circ$: adversarial floating-point searches found star positivity failing at minimum angle $15^\circ$ and never at $20^\circ$.

\subsubsection{Necessity from a locally quasi-uniform region, with explicit constants}

The failure of star witnesses shows that necessity must come from wider patches. The witness below is the curl of the Argyris interpolant of a fixed profile of width $Y$; on the exact half-plane its Rayleigh quotient is $Y(2B-A)/C=15.487$. Here $\theta$ is the minimum angle of $\Th$, and for $v\in\Zh$
\[
N_h(v,v)=\frac{C(v)}{Y}\Bigl[\frac{\mu Y}{h}-\mathcal Q(v)\Bigr],\qquad \mathcal Q(v):=\frac{Y\,(2B(v)-A(v))}{C(v)} .
\]

\paragraph{Profile.} $g(t)=t(1-t)^6$ for $-\tfrac18\le t\le1$, $g=0$ for $t\ge1$; $\phi(s)=(1-s^2)^6$ on $[-1,1]$, $0$ outside.
Both are $C^5$ with piecewise-polynomial, bounded sixth derivative. With $X=2Y$,
\[F(x,y)=Y\,\phi(x/X)\,g(y/Y),\qquad w=\curl F=(F_y,-F_x),\]
in coordinates in which $e$ lies on $\{y=0\}$, $x_0=0$, and the fluid is locally $\{y>0\}$. On the part of $\Oh$ below
$y=-Y/8$ (the far side of $P_h$) set $F:=0$. That part and the part near $x_0$ are separated by $P_h$, so no triangle meets
both (Lemma~\ref{pb:lem:geom}).

\paragraph{Scaled derivative bounds} (\texttt{notes/v6/consts/argyris\_consts.py}, exact critical points). Let
$m_j:=Y^{j-1}\sup|D^jF|$ over $\{y\ge-Y/8\}$. Bounding each partial derivative termwise, with
$\max_{|u|=1}|u_1|^i|u_2|^{j-i}$ in place of the monomials,
\[m_1\le3.6635,\quad m_2\le32.187,\quad m_3\le214.88,\quad m_6\le18918.2 .\]
\paragraph{Half-plane values} (exact, SymPy; $Y=1$):
$\hat A=\tfrac{4717730594816}{558289907175}=8.45032$, $\hat B=\tfrac{201326592}{16900975}=11.91213$,
$\hat C=\tfrac{16777216}{16900975}=0.992677$. So $\hat{\mathcal Q}=(2\hat B-\hat A)/\hat C=15.48734$, which agrees with the independent
value $15.487$ of \texttt{notes/v6/penalty/continuum\_check.py}.


\begin{lemma}[Explicit local bound]\label{pb:lem:arg}
Let $T$ be a triangle with minimum angle $\ge\theta$ and diameter $h_T$, and let $f\in C^5(T)$ have $D^5f$ Lipschitz, with
$\operatorname{ess\,sup}_T|D^6f|\le M_6$. Let $I$ be the Argyris ($C^1$-$P_5$) interpolant. Its DOFs are the values,
gradients and Hessians at the vertices and the unit-normal derivatives at the edge midpoints. Then on $T$, for $j=1,2$,
\[|D^j(f-If)|\le E_j(\theta)\,M_6\,h_T^{6-j},\]
where
\[E_j(\theta)=\frac1{(6-j)!}+\Bigl(\frac{\sqrt2}{\sin\theta\,s(\theta)}\Bigr)^{j}\Bigl[P_j+\bigl(\tfrac1{60}+2c_H\bigr)N_{0,j}
+\bigl(\tfrac1{3840}+\tfrac{\sqrt2\,c_H}{2\sin\theta}\bigr)(N_{1,j}+N_{2,j})\Bigr].\]
Here $s(\theta)=\min(\sqrt3/2,\sin2\theta)$ and $c_H=\tfrac{29}{1920}$. The reference constants are
\begin{gather*}P_1\le0.19597,\qquad P_2\le1.40308,\\ (N_{0,1},N_{1,1},N_{2,1})\le(1.8857,3.7713,3.7713),\qquad
(N_{0,2},N_{1,2},N_{2,2})\le(16.866,33.731,33.731).\end{gather*}
Numerically: $E_1=17.9,6.54,3.37,1.47,0.933,0.724$ and $E_2=3606,597,179,39.0,17.3,10.9$ for
$\theta=10^\circ,15^\circ,20^\circ,30^\circ,45^\circ,60^\circ$.
\end{lemma}

\begin{proof}
\begin{enumerate}[label=\arabic*.]
\item \emph{Taylor.} Let $a_0$ be the vertex with the largest angle $\alpha\in[\pi/3,\pi-2\theta]$, so
$\sin\alpha\ge s(\theta)$. Let $p$ be the degree-5 Taylor polynomial of $f$ at $a_0$ and $r=f-p$. Since $Ip=p$,
$f-If=r-Ir$. Along each segment $[a_0,x]\subset T$, $D^5f$ is Lipschitz, so the integral remainder gives
\[|D^mr(x)|\le M_6|x-a_0|^{6-m}/(6-m)!\le M_6h_T^{6-m}/(6-m)!,\]
and all derivatives of $r$ of order $\le5$ vanish at $a_0$.
\item \emph{Reference map.} Let $\Phi(\hat x)=J\hat x+a_0$ map $\hat T=\mathrm{conv}\{(0,0),(1,0),(0,1)\}$ onto $T$. The columns
of $J$ are the edges $E_1,E_2$ at $a_0$, of lengths $\ell_1,\ell_2$. By the law of sines,
$h_T\sin\theta\le\ell_i\le h_T$, and the longest edge $E_0$ (opposite $a_0$) has length $h_T$. Hence
\[\sigma_{\max}(J)\le\sqrt2h_T,\qquad \norm{J^{-1}}=\sigma_{\max}/\det J\le\frac{\sqrt{\ell_1^{-2}+\ell_2^{-2}}}{\sin\alpha}\le\frac{\sqrt2}{h_T\sin\theta\,s(\theta)} .\]
\item \emph{Pulled-back interpolant.} Let $\hat q=(Ir)\circ\Phi\in P_5$. Its vertex values, gradients and Hessians are those of
$\hat r=r\circ\Phi$, because these DOF sets are affine-equivalent. Only the edge DOFs differ. On reference edge $\hat e$ with
edge vector $\hat t$ (the reference DOF being $\nabla\hat q(\hat m)\cdot\hat\nu$, with
$\hat\nu\in\{(1,1),(-1,0),(0,-1)\}$), write $J^{-1}n_E=a\hat\nu+b\hat t$. Then
$\gamma_E:=\nabla\hat q(\hat m)\cdot\hat\nu=[n_E\cdot\nabla r(m_E)-b(Hf_E)'(\tfrac12)]/a$. Here $f_E(s)=\hat r(\hat p+s\hat t)$, and
$Hf_E$ is its quintic Hermite interpolant from $f,f',f''$ at $s=0,1$, which is the restriction of $\hat q$. Exactly,
$(Hf)'(\tfrac12)=\tfrac{15}8(f_1-f_0)-\tfrac7{16}(f_0'+f_1')+\tfrac1{32}(f_1''-f_0'')$.
\item \emph{Coefficient bounds.} Since $\det(J^{-1}n_E,\hat t)=\pm\ell_E/\det J$,
\[|a|^{-1}=\det J\,|\det(\hat\nu,\hat t)|/\ell_E,\qquad |b/a|\le|\hat\nu|\sigma_{\max}(J)/\ell_E .\]
This gives:
\begin{itemize}
\item on the hypotenuse ($E_0$): $|a|^{-1}\le2h_T$ and $|b/a|\le2$;
\item on the legs: $|a|^{-1}\le h_T$ and $|b/a|\le\sqrt2/\sin\theta$.
\end{itemize}
With $f_E^{(m)}=D^mr[E^m]$ and step 1:
\begin{itemize}
\item $|\gamma_{E_0}|\le M_6h_T^6(\tfrac2{120}+2c_H)$;
\item $|\gamma_{E_{1,2}}|\le M_6h_T^6(\tfrac{2^{-5}}{120}+\tfrac{\sqrt2}{\sin\theta}\tfrac{c_H}2)$, since the data at $a_0$ vanish.
\end{itemize}
The vertex DOFs at $a_1,a_2$ are $D^mr[E_{i_1},\dots,E_{i_m}]$, of size $\le M_6h_T^6/(6-m)!$.
\item \emph{Reference basis.} Expand $\hat q$ in the exact reference Argyris basis ($21\times21$ inverse over $\mathbb Q$).
Bound each basis function's partial derivatives of order $j$ on $\hat T$ by the largest absolute Bernstein coefficient, which
is a rigorous sup bound. Combine them into operator-norm bounds via the Frobenius norm. This gives
$|\hat D^j\hat q|\le M_6h_T^6[P_j+\dots]$, the bracket in $E_j$.
\item \emph{Back to $T$.} $|D^j(Ir)|\le\norm{J^{-1}}^j|\hat D^j\hat q|$. Add $|D^jr|$ from step 1.\qedhere
\end{enumerate}
\end{proof}
\emph{Check} (\texttt{notes/v6/consts/validate\_interp.py}, (NUMERICAL)). Physical Argyris interpolation of $F$ on 600 random triangles
($\theta\in[10^\circ,53^\circ]$, $h_T\in[0.01,0.2]$, random positions in $\operatorname{supp}F$): the ratio of actual error to
the bound is at most $1.8\cdot10^{-5}$ ($j=1$) and $1.4\cdot10^{-5}$ ($j=2$). The bound holds, and is very loose.


\begin{lemma}[Local geometry of $\Gh$]\label{pb:lem:geom}
Suppose every boundary edge meeting $B(x_0,3Y)$ has length $\le\kappa Y$, with $\kappa\le1$ and $Y\le10^{-2}$. Put
$X'=(2+\kappa)Y$, $\omega_X=\arcsin\!\bigl(X'/(1-\kappa^2Y^2/4)\bigr)$,
\[\delta=\tfrac14\kappa^2Y^2+\tfrac12\omega_X^2,\qquad \sigma=\tan\bigl(\omega_X+\arcsin(\kappa Y/2)\bigr).\]
Then on $|x|\le X'$ the part of $\Gh$ near $x_0$ is a graph $y=\varpi(x)$ with $-\delta\le\varpi\le0$ and $|\varpi'|\le\sigma$,
and $\Oh\cap\{|x|\le X',\,y>-Y/8\}=\{y>\varpi(x)\}$ (this uses $\delta\le Y/8$, true for $Y\le10^{-2}$).
\end{lemma}
\begin{proof}
\begin{itemize}
\item \emph{$\varpi\le0$.} $P_h$ is convex and $e$ is one of its edges, so $P_h$ lies in $\{y\le0\}$.
\item \emph{Polar description.} Let $O=(0,-d)$ be the centre. A boundary point on a chord of length $\ell'\le\kappa Y$ is
$O+r(\sin\omega,\cos\omega)$, with $r\ge\sqrt{1-\ell'^2/4}\ge1-\ell'^2/4$.
\item \emph{Depth.} So $y\ge-\ell'^2/4-(1-\cos\omega)$, while $|x|\le X'$ forces $|\omega|\le\omega_X$.
\item \emph{Slope.} A chord between polar angles $\omega_1,\omega_2$ has direction angle $(\omega_1+\omega_2)/2$, and
$|\omega_2-\omega_1|\le2\arcsin(\ell'/2)$.
\item \emph{Fluid side.} The region between $\varpi$ and $y=-Y/8$ lies in $P_h$ by convexity.\qedhere
\end{itemize}
\end{proof}

\begin{theorem}[Theorem A with explicit constants]\label{pb:thmA}
Let $e\in\EG$ have midpoint $x_0$, let $0<Y\le Y_0:=10^{-6}$ with $4Y<L-1$, let $0<\kappa\le\kappa_0(\theta)$ (table below), and suppose every $T\in\Th$ meeting $B(x_0,3Y)$ has $\diam T\le\kappa Y$. Let $v:=\curl I F$, with $I$ the global
Argyris interpolant. Then $v\in\Zh$ for every $k\ge4$ and $\mathcal Q(v)\ge15$. Hence $N_h$ is indefinite on $\Zh$ whenever
$\mu<15h/Y$. The values of $\kappa_0(\theta)$ are
\[\begin{array}{c|cccccc}\theta&10^\circ&15^\circ&20^\circ&30^\circ&45^\circ&60^\circ\\\hline
\kappa_0(\theta)&0.00379&0.00594&0.00803&0.0117&0.0143&0.0161\end{array}\]
(bisection roots of $\mathcal Q_{\rm lb}=15$ below, rounded down; the constants are computed in exact rational arithmetic and the final assembly in floating point, with rounding far below the stated digits).
\end{theorem}
\begin{proof}
\begin{enumerate}[label=\arabic*.]
\item \emph{Membership in $\Zh$.}
 \begin{itemize}
 \item $IF$ is globally $C^1$ and piecewise $P_5$, so $v$ is continuous, piecewise $P_4\subset P_k$, and pointwise
 divergence-free.
 \item $IF|_T=0$ on every $T$ not meeting $\operatorname{supp}F\cap\overline\Oh\subset\{|x|\le2Y,\ -\delta\le y\le Y\}\subset B(x_0,\sqrt5Y+\delta)$.
 \item So $\operatorname{supp}v\subset S:=\bigcup\{T:T\cap\operatorname{supp}F\ne\emptyset\}\subset\{|x|\le X',\ \varpi\le y\le(1+\kappa)Y\}$,
 which misses $\partial R$.
 \item Every $T\subset S$ lies in $\{y\ge-\delta\}$, where $F$ is given by the profile formula, so Lemma~\ref{pb:lem:arg} applies
 with $M_6=m_6Y^{-5}$ and $h_T\le\kappa Y$.
 \end{itemize}
 Write $e_F=F-IF$, $|S|\le2X'((1+\kappa)Y+\delta)$ and $\Lambda:=|\Gh\cap S|\le2X'\sqrt{1+\sigma^2}$.
\item \emph{Continuum field on the true geometry.} By Lemma~\ref{pb:lem:geom}, $\Oh\supset\{0<y<Y,\ |x|<X\}$, so
$A_{\Oh}(w)=\hat A+A_{\rm sl}$ with $0\le A_{\rm sl}\le4\sup|D^2F|^2\cdot2X\delta=16m_2^2\delta/Y$.
 \begin{itemize}
 \item \emph{Boundary term.} Using $n\,ds=(\varpi',-1)\,dx$ and moving from $y=0$ to $y=\varpi$:
 $|B_{\Gh}(w)-\hat B|\le\varepsilon_B:=4[\sigma m_1m_2+(\delta/Y)(m_1m_3+m_2^2)]$.
 \item \emph{Trace term.} $|C_{\Gh}(w)-Y\hat C|\le\varepsilon_C:=4Y[m_1^2(\sqrt{1+\sigma^2}-1)+2\sqrt{1+\sigma^2}\,m_1m_2\delta/Y]$.
 \end{itemize}
\item \emph{Interpolation.} $v-w=\curl e_F$, so $|\nabla(v-w)|_F=|D^2e_F|_F\le\sqrt2|D^2e_F|$ and $|v-w|=|\nabla e_F|$. By
Lemma~\ref{pb:lem:arg}, $|D^2e_F|\le E_2m_6\kappa^4/Y$ and $|\nabla e_F|\le E_1m_6\kappa^5$ on $S$. With $A(z)^{1/2}\le\sqrt2\norm{\nabla z}$:
\begin{gather*}A(v)^{1/2}\le(\hat A+A_{\rm sl})^{1/2}+\eta_A,\quad \eta_A=2E_2m_6\kappa^4\sqrt{|S|}/Y;\\
C(v)^{1/2}\le(Y\hat C+\varepsilon_C)^{1/2}+\eta_C,\quad \eta_C=E_1m_6\kappa^5\sqrt\Lambda;\end{gather*}
\[|B(v)-B_{\Gh}(w)|\le\varepsilon_B':=\frac{E_2m_6\kappa^4\sqrt\Lambda}{Y}\bigl((Y\hat C+\varepsilon_C)^{1/2}+\eta_C\bigr)+\frac{m_2\sqrt\Lambda}{Y}\eta_C,\]
using $B(v)-B(w)=\ip{\dn(v-w)}v+\ip{\dn w}{v-w}$.
\item \emph{Assembly.}
\[\mathcal Q(v)\ \ge\ \mathcal Q_{\rm lb}:=\frac{Y\bigl[2(\hat B-\varepsilon_B-\varepsilon_B')-((\hat A+A_{\rm sl})^{1/2}+\eta_A)^2\bigr]}{((Y\hat C+\varepsilon_C)^{1/2}+\eta_C)^2}.\]
Every error term is nondecreasing in $\kappa$ and in $Y$. $\mathcal Q_{\rm lb}$ is dimensionless: $|S|/Y^2$, $\Lambda/Y$,
$\delta/Y$ and $\sigma$ depend only on $\kappa$ and $Y$ through the formulas above. The table value $\kappa_0(\theta)$ is the
bisection root of $\mathcal Q_{\rm lb}(\kappa,\theta,Y_0)=15$, rounded down (\texttt{notes/v6/consts/assemble.py}, \texttt{assemble.log}). A
grid check below $(\kappa_0,Y_0)$ confirms $\min\mathcal Q_{\rm lb}\ge15.0000$. As $\kappa,Y\to0$, $\mathcal Q_{\rm lb}\to15.48734$.\qedhere
\end{enumerate}
\end{proof}

\begin{corollary}[$\rho$-scaling]\label{pb:corA}
Assume (M0). Suppose that for some edge $e\in\EG$, with midpoint $x_0$, every triangle meeting $B(x_0,3C_q|e|/\kappa_0)$ has
diameter $\le C_q|e|$, where $C_q\ge1$ and $C_q|e|/\kappa_0\le Y_0$. Then
\[\mu^\ast\ \ge\ \frac{15\kappa_0(\theta)}{C_q}\,\frac h{|e|}\ \ge\ \frac{15\,c_0\,\kappa_0(\theta)}{C_q}\,\rho ,\]
and the factor $c_0$ can be dropped when $e$ is a shortest edge of $\Gh$.
\end{corollary}
\begin{proof}
Take $Y=C_q|e|/\kappa_0$ in Theorem~\ref{pb:thmA}. For the second inequality, $|e|\le\hG\le\min|e'|/c_0$ by (M0).
\end{proof}
The smallness $Y_0=10^{-6}$ is crude but explicit: it forces $|e|\lesssim\kappa_0Y_0/C_q\approx10^{-8}$, so the corollary is a qualitative statement about the scaling $\mu^\ast\gtrsim\rho$, not a practical threshold. At $Y_0=3\cdot10^{-6}$ the values of $\kappa_0$ drop by about $15\%$; for $Y\gtrsim5\cdot10^{-6}$ the bound falls below $15$ even as $\kappa\to0$. The smallness comes almost entirely from the sliver and curvature terms $A_{\rm sl}$ and $\varepsilon_B$, which are $O(Y)$ with large profile constants.

Theorem~\ref{pb:thmA} needs no inf-sup constant, no compactness and none of (M1), (M1$'$), (M2); the only analytic inputs are Taylor's formula, the exact reference Argyris basis and elementary triangle geometry. For Clough--Tocher refinements the same proof applies with the HCT reference basis (\texttt{notes/v6/penalty}), but we have not computed those constants.

\emph{How far the explicit $\kappa_0$ is from the truth.} (NUMERICAL) (\texttt{realistic\_kappa.py}, \texttt{.log}). On half-plane meshes (squares of side $s$ split by a
diagonal, optionally with random vertex perturbation $\pm s/4$), $\mathcal Q(\curl IF)$ with $Y=1$ is:
\begin{center}\small\begin{tabular}{lccccccc}
\toprule
$\kappa=\diam T/Y$ & 1.41 & 1.13 & 0.85 & 0.71 & 0.42 & 0.14 & 0.07\\
\midrule
uniform ($45^\circ$) & 12.15 & 14.06 & 15.06 & 15.28 & 15.46 & 15.487 & 15.487\\
\bottomrule
\end{tabular}\end{center}
The perturbed meshes (min angle $9$--$27^\circ$) give $14.79$ at $\kappa=1.06$ and $15.28$ at $\kappa=0.88$. So the true
threshold is $\kappa\approx0.85$--$1$, against the proved $0.004$--$0.016$. The loss comes from the Taylor/sup-norm bound
($\approx7\cdot10^4$ in $E_j$, entering as a fourth root) and from the $\sin\theta$ factors.


\emph{Graded meshes.} The hypothesis of Theorem~\ref{pb:thmA} is not implied by (M0)--(M2): shape regularity only gives $\diam T\le C(\theta)(|e|+\dist(T,x_0))$ with $C(\theta)\ge1$, and this growth is attained.

\begin{proposition}[Obstruction; half-plane model]\label{pb:obst}
Take the dyadic boundary layer:
\begin{itemize}
\item row 0: unit squares on $y\in[0,1]$, split by the forward diagonal;
\item row $j\ge1$: squares of side $2^j$ on $y\in[2^j-1,2^{j+1}-1]$, each over two row-$(j-1)$ squares, cut into three
triangles from the hanging bottom midpoint.
\end{itemize}
This mesh has minimum angle $\arctan\frac12=26.57^\circ$, equal boundary edges, boundary stars $\cong\cS_3$, and
$\Theta\ge\sin45^\circ$ at every vertex; the worst vertex is a coarse-square corner, with consecutive angles $90^\circ$ and
$45^\circ$. For every $x_0$ on the boundary and every $Y>0$, $B(x_0,3Y)$ contains a triangle of diameter $\ge Y$. So the hypothesis of Theorem~\ref{pb:thmA} fails for every $\kappa<1$.
\end{proposition}
\begin{proof}
\begin{itemize}
\item \emph{$Y<1$.} The boundary triangle at $x_0$ has diameter $\ge1$.
\item \emph{$Y\ge1$.} The point $x_0+(0,2Y)$ lies in row $j$ with $2^j-1\le2Y<2^{j+1}-1$. Its triangle has diameter
$\ge2^j>(2Y+1)/2\ge Y$.\qedhere
\end{itemize}
\end{proof}

\paragraph{The witness still works on graded meshes} (NUMERICAL) (\texttt{graded\_check.py},
\texttt{graded\_check3.py}). The boundary edge is $1$, and "max $\kappa$" is the largest $\diam T/Y$ over the support.
\begin{center}\footnotesize\begin{tabular}{lcccccccc}
\toprule
$Y/\hG$ & 1 & 1.5 & 2 & 3 & 4 & 8 & 16 & 32\\
\midrule
dyadic (ratio 2) & 12.36 & 14.82 & 15.28 & 15.44 & 15.46 & 15.485 & 15.488 & 15.488\\
\quad max $\kappa$ & 1.41 & 1.49 & 1.12 & 0.75 & 1.12 & 1.12 & 1.12 & 1.12\\
uniform & 12.36 & 14.82 & 15.28 & 15.44 & 15.47 & 15.486 & 15.487 & 15.487\\
\midrule
triadic (ratio 3, $15.3^\circ$), $Y=1,2,3,9,27$ & \multicolumn{8}{l}{$12.36,\ 15.24,\ 15.42,\ 15.49,\ 15.48$ (max $\kappa$ up to $1.8$)}\\
\bottomrule
\end{tabular}\end{center}
On the dyadic layer, $\mathcal Q$ tends to the continuum value \emph{although} $\kappa\approx1.1$ at every scale. The
witness's quality is set by the fine cells near $\Gh$. The coarse cells carry little energy, and the true interpolation error
there is far below the worst-case bound. So on graded meshes the conclusion $\mu^\ast\gtrsim h/\hG$ appears to survive.
Proving it this way needs a \emph{sharp} interpolation bound on coarse cells, with $\kappa\sim1$, which a Taylor/sup-norm
argument cannot deliver. Two routes would close the general case:
\begin{enumerate}[label=(\roman*)]
\item a computer-assisted, shape-uniform bound of the Argyris error functional for this specific $F$;
\item the half-plane-limit positivity of Theorem~B in \texttt{notes/v6/penalty}, restricted to the explicit family
$\{\curl I_AF_Y\}$, which is compact under (M0).
\end{enumerate}
Neither is done. (open).


\emph{What is gained unconditionally.} $\mu^\ast\ge15\,h/Y_\ast$, where $Y_\ast$ is the smallest $Y\le Y_0$ such that some ball $B(x_0,3Y)$ centred at an edge midpoint meets only triangles of diameter $\le\kappa_0Y$. Since the boundary triangle has diameter $\ge c_0\hG$, $\kappa_0Y_\ast\ge c_0\hG$; equality up to constants holds if and only if about $3/\kappa_0$ quasi-uniform layers exist somewhere along $\Gh$. A uniform necessity statement under (M0)--(M2) alone remains open. A compactness reduction to positivity of the witness family on all half-plane limit meshes is possible (\texttt{notes/v6/penalty}, Theorem~B there), but it is not written out here and does not settle positivity.

\subsubsection{Why the certified constant is below the global threshold}

The global threshold is $\gamma^\ast=\hG\sup_{\Zh}(2B-A)/C$, where $A$ is taken over $\Oh$, $B,C$ over $\Gh$, and
$\hG=\max|e|$. Every patch witness gives a lower bound. For the Section~\ref{sec:numerics} meshes we computed the patch
thresholds with the penalty code of Section~\ref{sec:numerics} (bisection on the inertia of $A_\mu+rB^{\mathsf T}M^{-1}B$, fields vanishing
outside the patch). We maximised over all boundary triangles and single-vertex stars. \emph{win}$W$ is the first layer
restricted to $W$ columns centred at $\theta=0$, and the column $J$ is the patch made of the first $J$ full boundary layers, computed for $J\le2$ here (the log also has $J=3,4$, which change the values by at most $0.002$). The \emph{global} column is Table~\ref{tab:thresh}(a); for $N\le16$ the full mesh is $J=N/4\le4$ layers and the log reproduces it.

\begin{center}\small
\begin{tabular}{rccccccccc}
\toprule
$N$ & triangle & star & win2 & win4 & win8 & win16 & $J=1$ & $J=2$ & global\\
\midrule
8 & 6.13 & 13.71 & 13.91 & 14.65 & -- & -- & 14.68 & 14.74 & 14.74\\
16 & 8.93 & 16.58 & 16.89 & 17.91 & 17.94 & -- & 17.95 & 18.09 & 18.09\\
32 & 10.08 & 17.48 & 17.82 & 19.54 & 19.69 & 19.69 & 19.69 & 19.85 & 19.85\\
64 & 10.49 & 17.70 & 18.03 & 20.29 & 20.68 & 20.69 & 20.69 & 20.80 & 20.80\\
\bottomrule
\end{tabular}
\end{center}
\cert{code/lamstar\_family.py mesh}{logs/lamstar\_family\_mesh.log}

\emph{The limit geometry.} The meshes of Section~\ref{sec:numerics} do not have equal chords. The square points are
equispaced and projected radially, so $\max|e|/\min|e|$ grows from $1.43$ ($N=16$) to $1.97$ ($N=256$). The longest chords
are at the side midpoints $\theta=0,\pi/2,\ldots$. There the first layer, scaled by the chord, tends to columns
$[i,i+1]\times[0,\frac34]$ (layer height $(L-1)/2=\frac34$) split on the forward diagonal. The single-vertex star of that
geometry is $\cS_{3/4}$, with rim $(1,0),(1,\frac34),(0,\frac34),(-1,0)$. The general exact code (continuity, Dirichlet and
divergence constraints over $\mathbb Q$, exact barycentric moments) gives:

\begin{center}\small
\begin{tabular}{lcccccc}
\toprule
 & star & $W=2$ & $W=4$ & $W=8$ & $W=16$ & $W=32$\\
\midrule
1 layer & 17.7064 & 17.9759 & 20.4872 & 21.1726 & 21.3358 & 21.3758\\
2 layers & -- & 18.1738 & 20.6781 & 21.2432 & 21.3675 & 21.3974\\
3 layers & -- & 18.1916 & 20.6811 & 21.2462 & 21.3708 & --\\
\bottomrule
\end{tabular}
\end{center}

\emph{Exact certificates in the limit geometry:} $2B-A-17C>0$ on $\cS_{3/4}$ (dimension $14$, witness quotient $17.70640$; also
confirmed by \texttt{code/lamstar\_exact.py}), and $2B-A-20C>0$ on the $W=4$ window (dimension $41$, witness quotient
$20.48725$).
\cert{code/lamstar\_family.py limit; lamstar\_exact.run on $\cS_{3/4}$}{logs/lamstar\_family\_limit.log, logs/lamstar\_family\_limitstar.log}

The perturbation argument of Proposition~\ref{prop:penalty}(b) transfers these strict inequalities to the actual meshes near
$\theta=0$ for all sufficiently large $N$ (non-explicit $N_0$). All patch vertices keep at least three distinct edge lines, or
are Dirichlet corners where $v$ and $\nabla v$ vanish, so the kernel dimension is locally constant. This gives
$\gamma^\ast(N)>20$ for $N\ge N_0$. The finite-$N$ values in the first table agree: win4 is $20.29$ at $N=64$.

The gap between the certified $9$ and the measured $18$--$21$ is therefore accounted for, in order of size, by:
\begin{enumerate}
\item \emph{Shape.} The reference stars $\cS_3,\cS_4'$ have layer height $1$ per unit chord. The meshes have height $\frac34$ at the
longest chords, which is where $\gamma^\ast$ is decided because $\gamma$ is normalised by $\max|e|$. $\lambda^\ast$ grows fast as the
layer flattens: the exact witness quotients, which are lower bounds for $\lambda^\ast$, are $9.98$ for $\cS_3$ against $17.71$ for $\cS_{3/4}$.
\item \emph{Collective boundary modes.} The optimal witness spreads over $4$--$8$ boundary edges and adds about $3.7$
($17.7\to21.4$). Windows wider than about $16$ edges add nothing.
\item \emph{Depth.} In the limit geometry the second layer adds $0.20$ ($W=2$) and $0.19$ ($W=4$), decreasing to $0.02$ at $W=32$; the third layer adds $0.018$ at $W=2$ and at most $0.004$ for $W\ge4$. On the actual meshes the second layer adds $0.11$--$0.16$ (compare $J=1$ and $J=2$). $\gamma^\ast$ is a boundary-layer quantity
(compare $J=2$ with the global value).
\item \emph{Finite $N$.} Away from the limit geometry (curvature, the varying chord and height along the layer) $\gamma^\ast(N)$ is
lower. Richardson extrapolation of the window sequences (differences shrink by a factor of about $4$ per doubling of $W$) gives
$\gamma^\ast_\infty\approx21.39$ (one layer) and $21.41$ (two layers), so we expect $\gamma^\ast(N)\to\approx21.4$ on this mesh family.
\end{enumerate}
